% Seminar 3 - Pagina 3. Fragment pentru
% % Seminar 3 - Pagina 3. Fragment pentru
% % Seminar 3 - Pagina 3. Fragment pentru
% % Seminar 3 - Pagina 3. Fragment pentru
% \input{seminars/pages/seminar3-pagina3-ro}; fără preambul sau pagină nouă.
% NOTE EDITORIALE:
% - Manuscrisul definește o mulțime "numărabilă" printr-o bijecție cu N,
%   adică în sensul "infinit numărabil". Urmează observația despre
%   submulțimi: finite sau numărabile.
% - Simbolurile de lângă observația despre submulțimi sunt șterse/slabe;
%   lectura "finită sau numărabilă" este provizorie.
% - Exemplul Q ~ Z x Z este păstrat fără reconstruirea unei bijecții.
% - Indicii cifrelor zecimale y_{n1}, y_{n2} și z_0,z_1 sunt păstrați.
% - La contradicție, manuscrisul este șters parțial; z != y_n este
%   dedus din construcția imediat precedentă.
% - Manuscrisul nu tratează ambiguitatea reprezentărilor zecimale
%   (șirurile care se termină în 9). Nu am completat demonstrația.
% - Nota despre produsul aplicațiilor este distinctă de exercițiul 1.
% - Exercițiul 2 continuă pe pagina următoare a manuscrisului.

\subsection*{Seminar 3 - Pagina 3}

\noindent O mulțime $A$ se numește \emph{numărabilă} dacă există o
bijecție
\[
 f:\mathbb N\longrightarrow A.
\]
\noindent\textbf{Exemple.} $\mathbb N$, $\mathbb Z$,
$\mathbb Q\ (\sim\mathbb Z\times\mathbb Z)$.

\noindent Dacă $A$ este numărabilă și $B\subseteq A$, atunci $B$ este
finită sau numărabilă.

\medskip
\noindent\textbf{(1) $\mathbb R$ nu este numărabilă.}

\noindent\underline{Demonstrație (prin reducere la absurd).}
Presupunem că $\mathbb R$ este numărabilă. Atunci $[0,1)$ ar fi
numărabilă. Fie
\[
 [0,1)=\{y_n\mid n\in\mathbb N\},
\]
unde scriem dezvoltările zecimale
\[
\begin{aligned}
 y_0 &= 0,y_{01}y_{02}y_{03}\ldots,\\
 y_1 &= 0,y_{11}y_{12}y_{13}\ldots,\\
     &\ \vdots\\
 y_n &= 0,y_{n1}y_{n2}y_{n3}\ldots .
\end{aligned}
\]
Construim numărul
\[
 z:=0,z_0z_1z_2\ldots z_t\ldots\in[0,1)
\]
prin
\[
 z_0=\begin{cases}
       1+y_{01},&y_{01}\neq9,\\
       0,&y_{01}=9,
     \end{cases}
 \qquad
 z_t=\begin{cases}
       1+y_{t,t+1},&y_{t,t+1}\neq9,\\
       0,&y_{t,t+1}=9.
     \end{cases}
\]
Astfel $z\neq y_n$, pentru orice $n\in\mathbb N$, contradicție cu
$[0,1)=\{y_n\mid n\in\mathbb N\}$.
Deci $[0,1)$ nu este numărabilă, prin urmare nici
$\mathbb R$ nu este numărabilă.

\medskip
\noindent\emph{Notă despre produsul mulțimilor numărabile.}
Fie
\[
 A_1=\{x_1,x_2,\ldots,x_n,\ldots\},\qquad
 B_1=\{y_1,y_2,\ldots,y_n,\ldots\},
\]
cu aplicațiile $f_1:A_1\to\mathbb N$ și $f_2:B_1\to\mathbb N$.
În manuscris apare aplicația produs
\[
\begin{aligned}
 f_1\times f_2 &: A_1\times B_1\longrightarrow
                    \mathbb N\times\mathbb N,\\
 (f_1\times f_2)(x_i,y_j)
   &=\bigl(f_1(x_i),f_2(y_j)\bigr).
\end{aligned}
\]

\medskip
\noindent\textbf{Exercițiul 2.} Nu există o aplicație
$f:\mathbb R\to\mathbb R$ astfel încât
\[
 |f(x)-f(y)|>1\qquad(\forall x\neq y).
\]

\noindent\underline{Demonstrație (prin reducere la absurd).}
Presupunem că există $f:\mathbb R\to\mathbb R$ cu
$|f(x)-f(y)|>1$ pentru $x\neq y$. Atunci
\[
 x\neq y\ \Longrightarrow\ f(x)\neq f(y),
\]
și deci $f$ este injectivă. Notăm
\[
 \operatorname{im}(f)=\{f(x)\mid x\in\mathbb R\}.
\]
Din injectivitate rezultă
$\operatorname{im}(f)\sim\mathbb R$, deci
$\operatorname{im}(f)$ nu este numărabilă.
Dar
\[
 \operatorname{im}(f)
  =\operatorname{im}(f)\cap\mathbb R
  =\operatorname{im}(f)\cap
     \left(\bigcup_{n\in\mathbb Z}[n,n+1)\right)
  =\ \cdots
\]
% Demonstrația continuă pe pagina următoare a manuscrisului.
; fără preambul sau pagină nouă.
% NOTE EDITORIALE:
% - Manuscrisul definește o mulțime "numărabilă" printr-o bijecție cu N,
%   adică în sensul "infinit numărabil". Urmează observația despre
%   submulțimi: finite sau numărabile.
% - Simbolurile de lângă observația despre submulțimi sunt șterse/slabe;
%   lectura "finită sau numărabilă" este provizorie.
% - Exemplul Q ~ Z x Z este păstrat fără reconstruirea unei bijecții.
% - Indicii cifrelor zecimale y_{n1}, y_{n2} și z_0,z_1 sunt păstrați.
% - La contradicție, manuscrisul este șters parțial; z != y_n este
%   dedus din construcția imediat precedentă.
% - Manuscrisul nu tratează ambiguitatea reprezentărilor zecimale
%   (șirurile care se termină în 9). Nu am completat demonstrația.
% - Nota despre produsul aplicațiilor este distinctă de exercițiul 1.
% - Exercițiul 2 continuă pe pagina următoare a manuscrisului.

\subsection*{Seminar 3 - Pagina 3}

\noindent O mulțime $A$ se numește \emph{numărabilă} dacă există o
bijecție
\[
 f:\mathbb N\longrightarrow A.
\]
\noindent\textbf{Exemple.} $\mathbb N$, $\mathbb Z$,
$\mathbb Q\ (\sim\mathbb Z\times\mathbb Z)$.

\noindent Dacă $A$ este numărabilă și $B\subseteq A$, atunci $B$ este
finită sau numărabilă.

\medskip
\noindent\textbf{(1) $\mathbb R$ nu este numărabilă.}

\noindent\underline{Demonstrație (prin reducere la absurd).}
Presupunem că $\mathbb R$ este numărabilă. Atunci $[0,1)$ ar fi
numărabilă. Fie
\[
 [0,1)=\{y_n\mid n\in\mathbb N\},
\]
unde scriem dezvoltările zecimale
\[
\begin{aligned}
 y_0 &= 0,y_{01}y_{02}y_{03}\ldots,\\
 y_1 &= 0,y_{11}y_{12}y_{13}\ldots,\\
     &\ \vdots\\
 y_n &= 0,y_{n1}y_{n2}y_{n3}\ldots .
\end{aligned}
\]
Construim numărul
\[
 z:=0,z_0z_1z_2\ldots z_t\ldots\in[0,1)
\]
prin
\[
 z_0=\begin{cases}
       1+y_{01},&y_{01}\neq9,\\
       0,&y_{01}=9,
     \end{cases}
 \qquad
 z_t=\begin{cases}
       1+y_{t,t+1},&y_{t,t+1}\neq9,\\
       0,&y_{t,t+1}=9.
     \end{cases}
\]
Astfel $z\neq y_n$, pentru orice $n\in\mathbb N$, contradicție cu
$[0,1)=\{y_n\mid n\in\mathbb N\}$.
Deci $[0,1)$ nu este numărabilă, prin urmare nici
$\mathbb R$ nu este numărabilă.

\medskip
\noindent\emph{Notă despre produsul mulțimilor numărabile.}
Fie
\[
 A_1=\{x_1,x_2,\ldots,x_n,\ldots\},\qquad
 B_1=\{y_1,y_2,\ldots,y_n,\ldots\},
\]
cu aplicațiile $f_1:A_1\to\mathbb N$ și $f_2:B_1\to\mathbb N$.
În manuscris apare aplicația produs
\[
\begin{aligned}
 f_1\times f_2 &: A_1\times B_1\longrightarrow
                    \mathbb N\times\mathbb N,\\
 (f_1\times f_2)(x_i,y_j)
   &=\bigl(f_1(x_i),f_2(y_j)\bigr).
\end{aligned}
\]

\medskip
\noindent\textbf{Exercițiul 2.} Nu există o aplicație
$f:\mathbb R\to\mathbb R$ astfel încât
\[
 |f(x)-f(y)|>1\qquad(\forall x\neq y).
\]

\noindent\underline{Demonstrație (prin reducere la absurd).}
Presupunem că există $f:\mathbb R\to\mathbb R$ cu
$|f(x)-f(y)|>1$ pentru $x\neq y$. Atunci
\[
 x\neq y\ \Longrightarrow\ f(x)\neq f(y),
\]
și deci $f$ este injectivă. Notăm
\[
 \operatorname{im}(f)=\{f(x)\mid x\in\mathbb R\}.
\]
Din injectivitate rezultă
$\operatorname{im}(f)\sim\mathbb R$, deci
$\operatorname{im}(f)$ nu este numărabilă.
Dar
\[
 \operatorname{im}(f)
  =\operatorname{im}(f)\cap\mathbb R
  =\operatorname{im}(f)\cap
     \left(\bigcup_{n\in\mathbb Z}[n,n+1)\right)
  =\ \cdots
\]
% Demonstrația continuă pe pagina următoare a manuscrisului.
; fără preambul sau pagină nouă.
% NOTE EDITORIALE:
% - Manuscrisul definește o mulțime "numărabilă" printr-o bijecție cu N,
%   adică în sensul "infinit numărabil". Urmează observația despre
%   submulțimi: finite sau numărabile.
% - Simbolurile de lângă observația despre submulțimi sunt șterse/slabe;
%   lectura "finită sau numărabilă" este provizorie.
% - Exemplul Q ~ Z x Z este păstrat fără reconstruirea unei bijecții.
% - Indicii cifrelor zecimale y_{n1}, y_{n2} și z_0,z_1 sunt păstrați.
% - La contradicție, manuscrisul este șters parțial; z != y_n este
%   dedus din construcția imediat precedentă.
% - Manuscrisul nu tratează ambiguitatea reprezentărilor zecimale
%   (șirurile care se termină în 9). Nu am completat demonstrația.
% - Nota despre produsul aplicațiilor este distinctă de exercițiul 1.
% - Exercițiul 2 continuă pe pagina următoare a manuscrisului.

\subsection*{Seminar 3 - Pagina 3}

\noindent O mulțime $A$ se numește \emph{numărabilă} dacă există o
bijecție
\[
 f:\mathbb N\longrightarrow A.
\]
\noindent\textbf{Exemple.} $\mathbb N$, $\mathbb Z$,
$\mathbb Q\ (\sim\mathbb Z\times\mathbb Z)$.

\noindent Dacă $A$ este numărabilă și $B\subseteq A$, atunci $B$ este
finită sau numărabilă.

\medskip
\noindent\textbf{(1) $\mathbb R$ nu este numărabilă.}

\noindent\underline{Demonstrație (prin reducere la absurd).}
Presupunem că $\mathbb R$ este numărabilă. Atunci $[0,1)$ ar fi
numărabilă. Fie
\[
 [0,1)=\{y_n\mid n\in\mathbb N\},
\]
unde scriem dezvoltările zecimale
\[
\begin{aligned}
 y_0 &= 0,y_{01}y_{02}y_{03}\ldots,\\
 y_1 &= 0,y_{11}y_{12}y_{13}\ldots,\\
     &\ \vdots\\
 y_n &= 0,y_{n1}y_{n2}y_{n3}\ldots .
\end{aligned}
\]
Construim numărul
\[
 z:=0,z_0z_1z_2\ldots z_t\ldots\in[0,1)
\]
prin
\[
 z_0=\begin{cases}
       1+y_{01},&y_{01}\neq9,\\
       0,&y_{01}=9,
     \end{cases}
 \qquad
 z_t=\begin{cases}
       1+y_{t,t+1},&y_{t,t+1}\neq9,\\
       0,&y_{t,t+1}=9.
     \end{cases}
\]
Astfel $z\neq y_n$, pentru orice $n\in\mathbb N$, contradicție cu
$[0,1)=\{y_n\mid n\in\mathbb N\}$.
Deci $[0,1)$ nu este numărabilă, prin urmare nici
$\mathbb R$ nu este numărabilă.

\medskip
\noindent\emph{Notă despre produsul mulțimilor numărabile.}
Fie
\[
 A_1=\{x_1,x_2,\ldots,x_n,\ldots\},\qquad
 B_1=\{y_1,y_2,\ldots,y_n,\ldots\},
\]
cu aplicațiile $f_1:A_1\to\mathbb N$ și $f_2:B_1\to\mathbb N$.
În manuscris apare aplicația produs
\[
\begin{aligned}
 f_1\times f_2 &: A_1\times B_1\longrightarrow
                    \mathbb N\times\mathbb N,\\
 (f_1\times f_2)(x_i,y_j)
   &=\bigl(f_1(x_i),f_2(y_j)\bigr).
\end{aligned}
\]

\medskip
\noindent\textbf{Exercițiul 2.} Nu există o aplicație
$f:\mathbb R\to\mathbb R$ astfel încât
\[
 |f(x)-f(y)|>1\qquad(\forall x\neq y).
\]

\noindent\underline{Demonstrație (prin reducere la absurd).}
Presupunem că există $f:\mathbb R\to\mathbb R$ cu
$|f(x)-f(y)|>1$ pentru $x\neq y$. Atunci
\[
 x\neq y\ \Longrightarrow\ f(x)\neq f(y),
\]
și deci $f$ este injectivă. Notăm
\[
 \operatorname{im}(f)=\{f(x)\mid x\in\mathbb R\}.
\]
Din injectivitate rezultă
$\operatorname{im}(f)\sim\mathbb R$, deci
$\operatorname{im}(f)$ nu este numărabilă.
Dar
\[
 \operatorname{im}(f)
  =\operatorname{im}(f)\cap\mathbb R
  =\operatorname{im}(f)\cap
     \left(\bigcup_{n\in\mathbb Z}[n,n+1)\right)
  =\ \cdots
\]
% Demonstrația continuă pe pagina următoare a manuscrisului.
; fără preambul sau pagină nouă.
% NOTE EDITORIALE:
% - Manuscrisul definește o mulțime "numărabilă" printr-o bijecție cu N,
%   adică în sensul "infinit numărabil". Urmează observația despre
%   submulțimi: finite sau numărabile.
% - Simbolurile de lângă observația despre submulțimi sunt șterse/slabe;
%   lectura "finită sau numărabilă" este provizorie.
% - Exemplul Q ~ Z x Z este păstrat fără reconstruirea unei bijecții.
% - Indicii cifrelor zecimale y_{n1}, y_{n2} și z_0,z_1 sunt păstrați.
% - La contradicție, manuscrisul este șters parțial; z != y_n este
%   dedus din construcția imediat precedentă.
% - Manuscrisul nu tratează ambiguitatea reprezentărilor zecimale
%   (șirurile care se termină în 9). Nu am completat demonstrația.
% - Nota despre produsul aplicațiilor este distinctă de exercițiul 1.
% - Exercițiul 2 continuă pe pagina următoare a manuscrisului.

\subsection*{Seminar 3 - Pagina 3}

\noindent O mulțime $A$ se numește \emph{numărabilă} dacă există o
bijecție
\[
 f:\mathbb N\longrightarrow A.
\]
\noindent\textbf{Exemple.} $\mathbb N$, $\mathbb Z$,
$\mathbb Q\ (\sim\mathbb Z\times\mathbb Z)$.

\noindent Dacă $A$ este numărabilă și $B\subseteq A$, atunci $B$ este
finită sau numărabilă.

\medskip
\noindent\textbf{(1) $\mathbb R$ nu este numărabilă.}

\noindent\underline{Demonstrație (prin reducere la absurd).}
Presupunem că $\mathbb R$ este numărabilă. Atunci $[0,1)$ ar fi
numărabilă. Fie
\[
 [0,1)=\{y_n\mid n\in\mathbb N\},
\]
unde scriem dezvoltările zecimale
\[
\begin{aligned}
 y_0 &= 0,y_{01}y_{02}y_{03}\ldots,\\
 y_1 &= 0,y_{11}y_{12}y_{13}\ldots,\\
     &\ \vdots\\
 y_n &= 0,y_{n1}y_{n2}y_{n3}\ldots .
\end{aligned}
\]
Construim numărul
\[
 z:=0,z_0z_1z_2\ldots z_t\ldots\in[0,1)
\]
prin
\[
 z_0=\begin{cases}
       1+y_{01},&y_{01}\neq9,\\
       0,&y_{01}=9,
     \end{cases}
 \qquad
 z_t=\begin{cases}
       1+y_{t,t+1},&y_{t,t+1}\neq9,\\
       0,&y_{t,t+1}=9.
     \end{cases}
\]
Astfel $z\neq y_n$, pentru orice $n\in\mathbb N$, contradicție cu
$[0,1)=\{y_n\mid n\in\mathbb N\}$.
Deci $[0,1)$ nu este numărabilă, prin urmare nici
$\mathbb R$ nu este numărabilă.

\medskip
\noindent\emph{Notă despre produsul mulțimilor numărabile.}
Fie
\[
 A_1=\{x_1,x_2,\ldots,x_n,\ldots\},\qquad
 B_1=\{y_1,y_2,\ldots,y_n,\ldots\},
\]
cu aplicațiile $f_1:A_1\to\mathbb N$ și $f_2:B_1\to\mathbb N$.
În manuscris apare aplicația produs
\[
\begin{aligned}
 f_1\times f_2 &: A_1\times B_1\longrightarrow
                    \mathbb N\times\mathbb N,\\
 (f_1\times f_2)(x_i,y_j)
   &=\bigl(f_1(x_i),f_2(y_j)\bigr).
\end{aligned}
\]

\medskip
\noindent\textbf{Exercițiul 2.} Nu există o aplicație
$f:\mathbb R\to\mathbb R$ astfel încât
\[
 |f(x)-f(y)|>1\qquad(\forall x\neq y).
\]

\noindent\underline{Demonstrație (prin reducere la absurd).}
Presupunem că există $f:\mathbb R\to\mathbb R$ cu
$|f(x)-f(y)|>1$ pentru $x\neq y$. Atunci
\[
 x\neq y\ \Longrightarrow\ f(x)\neq f(y),
\]
și deci $f$ este injectivă. Notăm
\[
 \operatorname{im}(f)=\{f(x)\mid x\in\mathbb R\}.
\]
Din injectivitate rezultă
$\operatorname{im}(f)\sim\mathbb R$, deci
$\operatorname{im}(f)$ nu este numărabilă.
Dar
\[
 \operatorname{im}(f)
  =\operatorname{im}(f)\cap\mathbb R
  =\operatorname{im}(f)\cap
     \left(\bigcup_{n\in\mathbb Z}[n,n+1)\right)
  =\ \cdots
\]
% Demonstrația continuă pe pagina următoare a manuscrisului.
